\subsection{Penerapan Nilai-Nilai Pancasila}
Nilai-nilai pancasila tidak hanya dihafalkan, ataupun disosialisasikan. Tetapi sungguh-sungguh dibudayakan.

\subsubsection{Strategi Perubahan dan Pengalaman}
\begin{itemize}
\item Melakukan perubahan dan pengalaman dan reaktualisasi pemahaman terhadap pancasila.
\item Mendorong wujudnya keadilan sosial.
\item Menguatkan internalisasi nilai-nilai pancasila ke dalam produk UUD.
\item Menumbuhkan keteladanan
\end{itemize}

\subsubsection{Pancasila Dapat Lestari dan Berkembang}
\begin{enumerate}
\item Dilingkungan keluarga: Kasih sayang
\item Dilingkungan sekolah: Toleransi terhadap keberagaman siswa
\item Dilingkungan masyarakat: Menghormati warga
\item Dilingkungan bangsa dan bernegara: Toleransi keberagaman
\end{enumerate}

\[
\text{Dinamika UUD 1945}
\begin{cases}
  \text{Konstitusi Tertulis } \rightarrow \text{UUD} \\
  \text{Konstitusi non Tertulis}
\end{cases}
\]

\textbf{Konstitusi: }Segala ketentuan dan aturan mengenai ketatanegaraan.


\textbf{UUD NRI: }Tidak dapat dipisahkan.


\noindent Menurut J.G Steenbeek $\rightarrow$ 3 pokok muatan konstitusi:

\begin{enumerate}
\item Pengaturan tentang pengaturan hak.
\item Pengaturan susunan ketatanegaraan.
\item Tugas-tugas dasar ketatanegaraan.
\end{enumerate}


\noindent \textbf{Pembukaan UUD: }Mempunyai kedudukan atas pasal-pasal UUD:
\begin{enumerate}
\item Alinea Pertama: Mengatur pengaturan tentang hak kodrat
\item Alinea Kedua: Merupakan konsekuensi logis dari pernyataan kemerdekaan.
\item Alinea Ketiga: Pengakuan religius dan nilai moral.
\item Alinea Keempat: Berisi tujuan umum dan khusus.
\end{enumerate}

\begin{enumerate}
\item \textbf{14 November: }Presidensial $\rightarrow$ parlementer.
\item \textbf{27 Desember 1945: }Konstitusi RIS.
\item \textbf{14 Agustus 1945: }UUDS 1950
\item \textbf{5 Juli 1945: }Dekrit Presiden \textit{``Konstitusi kembali ke UUD 1945 dan UUDS tidak berlaku lagi.''}
\end{enumerate}

\subsubsection{Definisi Kata}
\noindent \textbf{Serikat: }Negara yang disusun jamak/terdapat beberapa negara bagian.


\noindent \textbf{Republik: }Negara yang kepala negaranya disebut presiden yang dipilih rakyat.


\noindent \textbf{Presidensial: }Presiden sebagai kepala negara sekaligus juga kepada pemerintahan dibantu oleh mentri-mentri.


\noindent \textbf{Parlemen: }Memgang kekuasaan yang besar dalam mengontrol pemerintahan dan membuat kebijakan.
